% GENERATED FILE. Edit canonical lessons, then run npm run generate:books.
% canonical-source-hash: fnv1a64:8f04f3258bc7e3b9
% canonical-lessons: FR-C194-vallee, FR-C194-temperature, FR-C194-papillon, FR-C194-mouche, FR-C194-moustique

\chapter{Valley, Temperature, Butterfly, Fly, Mosquito}
\label{ch:fr-valley-temperature-butterfly-fly-mosquito}
\hlchaptermodality{194}

\begin{chapteropening}
\textbf{By the end of this chapter:} \emph{I can say a valley, temperature, a butterfly, a fly and a mosquito.}

Complete the last lesson of chapter 194: I can say a valley, temperature, a butterfly, a fly and a mosquito.
\end{chapteropening}

\section[la vallée]{la vallée — a valley}

\label{lesson:FR-C194-vallee}



\begin{quote}
Before the new one: say the French for a saucepan, then the French for a frying pan.
\end{quote}

\subsection*{You'll want to know: la vallée}
\textbf{la vallée} — \textquotedblleft{}a valley\textquotedblright{}.

\begin{cousinweb}
\textbf{vallée} — built in Old French from \emph{val}, \textquotedblleft{}a valley\textquotedblright{}, $\leftarrow$ Latin \emph{vallis}, with \emph{-ée}. (EtymDB 2.1)
\end{cousinweb}

\begin{grammarlens}[title={What you've built}]
One more word for this chapter.
\end{grammarlens}

\subsection*{Your turn}
\begin{itemize}
\raggedright
  \item \emph{Say it:} \emph{la vallée}
  \item \emph{Say it:} \emph{la vallée}, once more
  \item \emph{Say it:} say \emph{la vallée}
  \item \emph{Recall:} say the French for a saucepan, then the French for a frying pan, then say \emph{la vallée} again
\end{itemize}

\begin{tcolorbox}[breakable,colback=teal!4,colframe=teal!35!black,title={Before you move on}]
What does la vallée mean? (\textquotedblleft{}A valley\textquotedblright{}.) Say it once more.
\end{tcolorbox}
\section[la température]{la température — temperature}

\label{lesson:FR-C194-temperature}



\begin{quote}
Before the new one: say the French for a frying pan, then the French for a valley.
\end{quote}

\subsection*{You'll want to know: la température}
\textbf{la température} — \textquotedblleft{}temperature\textquotedblright{}.

\begin{cousinweb}
\textbf{température} $\leftarrow$ Latin \emph{temperatura}, \textquotedblleft{}due measure, a mixing\textquotedblright{}, from \emph{temperare}, \textquotedblleft{}to mix in proportion\textquotedblright{} — a learned borrowing. (EtymDB 2.1)
\end{cousinweb}

\begin{grammarlens}[title={What you've built}]
One more word for this chapter.
\end{grammarlens}

\subsection*{Your turn}
\begin{itemize}
\raggedright
  \item \emph{Say it:} \emph{la température}
  \item \emph{Say it:} \emph{la température}, once more
  \item \emph{Say it:} say \emph{la température}
  \item \emph{Recall:} say the French for a frying pan, then the French for a valley, then say \emph{la température} again
\end{itemize}

\begin{tcolorbox}[breakable,colback=teal!4,colframe=teal!35!black,title={Before you move on}]
What does la température mean? (\textquotedblleft{}Temperature\textquotedblright{}.) Say it once more.
\end{tcolorbox}
\section[le papillon]{le papillon — a butterfly}

\label{lesson:FR-C194-papillon}



\begin{quote}
Before the new one: say the French for a valley, then the French for temperature.
\end{quote}

\subsection*{You'll want to know: le papillon}
\textbf{le papillon} — \textquotedblleft{}a butterfly\textquotedblright{}.

\begin{cousinweb}
\textbf{papillon} $\leftarrow$ Latin \emph{papilio}, \emph{papilionis}, \textquotedblleft{}a butterfly\textquotedblright{} (EtymDB 2.1) — semi-learned: it kept its middle \emph{p}. The fully inherited twin, with \emph{v}, became the word for a tent, English \emph{pavilion}.
\end{cousinweb}

\begin{grammarlens}[title={What you've built}]
One more word for this chapter.
\end{grammarlens}

\subsection*{Your turn}
\begin{itemize}
\raggedright
  \item \emph{Say it:} \emph{le papillon}
  \item \emph{Say it:} \emph{le papillon}, once more
  \item \emph{Say it:} say \emph{le papillon}
  \item \emph{Recall:} say the French for a valley, then the French for temperature, then say \emph{le papillon} again
\end{itemize}

\begin{tcolorbox}[breakable,colback=teal!4,colframe=teal!35!black,title={Before you move on}]
What does le papillon mean? (\textquotedblleft{}A butterfly\textquotedblright{}.) Say it once more.
\end{tcolorbox}
\section[la mouche]{la mouche — a fly}

\label{lesson:FR-C194-mouche}



\begin{quote}
Before the new one: say the French for temperature, then the French for a butterfly.
\end{quote}

\subsection*{You'll want to know: la mouche}
\textbf{la mouche} — \textquotedblleft{}a fly\textquotedblright{}.

\begin{cousinweb}
\textbf{mouche} $\leftarrow$ Old French \emph{mousche} $\leftarrow$ Latin \emph{musca}, \textquotedblleft{}a fly\textquotedblright{} — inherited. (EtymDB 2.1)
\end{cousinweb}

\begin{grammarlens}[title={What you've built}]
One more word for this chapter.
\end{grammarlens}

\subsection*{Your turn}
\begin{itemize}
\raggedright
  \item \emph{Say it:} \emph{la mouche}
  \item \emph{Say it:} \emph{la mouche}, once more
  \item \emph{Say it:} say \emph{la mouche}
  \item \emph{Recall:} say the French for temperature, then the French for a butterfly, then say \emph{la mouche} again
\end{itemize}

\begin{tcolorbox}[breakable,colback=teal!4,colframe=teal!35!black,title={Before you move on}]
What does la mouche mean? (\textquotedblleft{}A fly\textquotedblright{}.) Say it once more.
\end{tcolorbox}
\section[le moustique]{le moustique — a mosquito}

\label{lesson:FR-C194-moustique}



\begin{quote}
Before the new one: say the French for a butterfly, then the French for a fly.
\end{quote}

\subsection*{You'll want to know: le moustique}
\textbf{le moustique} — \textquotedblleft{}a mosquito\textquotedblright{}.

\begin{cousinweb}
\textbf{moustique} $\leftarrow$ Spanish \emph{mosquito}, \textquotedblleft{}a little fly\textquotedblright{}, from \emph{mosca}, $\leftarrow$ Latin \emph{musca} — borrowed. (EtymDB 2.1)
\end{cousinweb}

\begin{grammarlens}[title={What you've built}]
One more word for this chapter.
\end{grammarlens}

\subsection*{Your turn}
\begin{itemize}
\raggedright
  \item \emph{Say it:} \emph{le moustique}
  \item \emph{Say it:} \emph{le moustique}, once more
  \item \emph{Say it:} say \emph{le moustique}
  \item \emph{Recall:} say the French for a butterfly, then the French for a fly, then say \emph{le moustique} again
\end{itemize}

\begin{tcolorbox}[breakable,colback=teal!4,colframe=teal!35!black,title={Before you move on}]
What does le moustique mean? (\textquotedblleft{}A mosquito\textquotedblright{}.) Say it once more.
\end{tcolorbox}
